% ============================================================
% introduccion.tex — Introducción y objetivos del trabajo
% ============================================================
\section{Introducción}

El fraude con tarjetas de crédito es un fenómeno delictivo que genera
pérdidas millonarias a nivel mundial. Según la Nilson Report, las pérdidas
globales por fraude en pagos con tarjeta superaron los 32 mil millones de
dólares en 2021, con proyecciones de crecimiento sostenido \cite{nilson2021}.
La detección temprana de estas actividades ilícitas es crucial para proteger
tanto a los consumidores como a las instituciones financieras.

Los métodos tradicionales de detección de fraude, basados en reglas
heurísticas, presentan limitaciones importantes frente a la evolución
constante de las tácticas fraudulentas. El aprendizaje automático ofrece
una alternativa adaptativa y escalable, capaz de identificar patrones
complejos y no lineales en grandes volúmenes de datos transaccionales
\cite{dal2015learned}.

El principal desafío técnico de este problema es el severo desbalance de
clases: en un entorno real, la proporción de transacciones fraudulentas
es menor al 1\% del total. Esto provoca que los clasificadores estándar
tiendan a predecir siempre la clase mayoritaria, logrando alta exactitud
pero fallando en detectar los fraudes, que son precisamente el objetivo
\cite{chawla2002smote}.

Este trabajo aborda el problema mediante la metodología CRISP-DM
\cite{wirth2000crisp}, aplicando técnicas de sobremuestreo sintético
(SMOTE) y comparando el rendimiento de cuatro algoritmos de clasificación.
El objetivo es construir un modelo robusto que maximice el recall de la
clase fraude, minimizando así las pérdidas financieras causadas por
transacciones no detectadas.

\subsection{Objetivos}
\begin{itemize}
  \item Implementar un pipeline completo de detección de fraude siguiendo CRISP-DM.
  \item Comparar el rendimiento de Regresión Logística, Random Forest,
        XGBoost y Gradient Boosting.
  \item Manejar el desbalance de clases mediante SMOTE.
  \item Desplegar el modelo como servicio web con visualización interactiva.
\end{itemize}
